% Three scalar price bounds used by the selected trace construction.
\section{Three local scalar price bounds}
\label{sec:ex-lattice-trace-medal-prices}

This section proves three scalar price conversions used by the selected
trace construction. The arguments are independent of the construction of the
local response and chart records that consume these bounds.

Fix a block side \(L\), and retain \(\beta(x)\), \(H(x)\), and
\(\chi(x)=H(x)^{1/2}\) from Section~\ref{sec:ex-lattice-trace-local-rate-selection}.
For a fixed price \(a>0\), write
\[
                 \mathsf s_a(x)=a/\chi(x).
\]
A negative strict-power rate has the following elementary consequence: if
\(f\in\mathcal R_\alpha\) with \(\alpha<0\), then
\(\chi(x)f(x)\to0\). Indeed \(\chi\in\mathcal R_0\), so the product has
rate \(\alpha\); choose \(v\in(\alpha,0)\), and use
\(\beta(x)^v\to0\).

\begin{lemma}[Cloud price]
\label{lem:ex-lattice-medal-cloud-price}
Let \(Q\) be a target with rate exponent \(q<0\). Suppose its poster
coefficient \(p_Q(x)\) has strict-power rate \(q/4\), and define the
cloud response for a cloud \(Z\) on a representation \(R\) by
\[
 \mathcal M_Q(Z;x)=m_R(Z)\exp(-\epsilon_Q N_R(Z))p_Q(x),
 \qquad m_R(Z)\ge0,\quad N_R(Z)\ge0,\quad \epsilon_Q>0.
\]
Then for every \(a>0\), eventually, uniformly over all clouds,
\[
             |\mathcal M_Q(Z;x)|\le \mathsf s_a(x)m_R(Z).
 \tag{1}\label{eq:ex-lattice-cloud-price}
\]
\end{lemma}
\begin{proof}
Since \(q<0\), choose \(v\) with \(q/4<v<0\). The rate hypothesis
implies \(|p_Q(x)|\le a/\chi(x)\) eventually: equivalently
\(\chi|p_Q|\to0\), and then make it at most \(a\). The attenuation obeys
\(0<\exp(-\epsilon_QN_R(Z))\le1\), independently of the cloud. Multiply
these inequalities by the nonnegative mass \(m_R(Z)\).
\end{proof}

\begin{lemma}[Forced gas price]
\label{lem:ex-lattice-medal-gas-price}
Let \(g(x)=\exp(-H(x)^{1/5}/2)\). For every \(a>0\), eventually
\[
                    \operatorname{DoorsRate}(L)g(x)\le\mathsf s_a(x).
 \tag{2}\label{eq:ex-lattice-gas-price}
\]
\end{lemma}
\begin{proof}
Set \(t=\log\log\beta(x)\), so \(\log H=t^2\) and
\(\log\beta=e^t\). Then
\[
\log\bigl(\chi(x)g(x)\bigr)=\tfrac12t^2-\tfrac12e^{t^2/5}\longrightarrow-\infty,
\]
because \(e^{t^2/5}\) grows faster than \(t^2\). Thus
\(\chi g\to0\); multiplying by the fixed positive \(\operatorname{DoorsRate}(L)\)
and making the product at most \(a\) proves (2).
\end{proof}

\begin{lemma}[Tuition price]
\label{lem:ex-lattice-medal-tuition-price}
For a target with goal \(\beta(x)^q\), define its tuition at fixed integers
\(s,n\ge0\) by
\[
 \mathcal T_Q(x;s,n)=A(x)^{sn}\,P_{1/2}(sn)\,\beta(x)^{-q}
                         \exp(-b(x)/2),
 \tag{3}\label{eq:ex-lattice-tuition}
\]
where \(A\) and the polynomial price \(P_{1/2}(sn)\) have rate zero and
\(b(x)=\lceil(\log x)^2\rceil\). Then
\(\mathcal T_Q(\cdot;s,n)\in\mathcal R_z\) for every real \(z\). In
particular, for every \(a>0\), eventually
\[
                       \mathcal T_Q(x;s,n)\le\mathsf s_a(x).
 \tag{4}\label{eq:ex-lattice-tuition-price}
\]
\end{lemma}
\begin{proof}
The first two factors in (3) have rate zero, and \(\beta^{-q}\) has rate
\(-q\). The last factor has every strict-power rate: for every fixed
\(c>0\), \(b(x)/2\ge c\log\beta(x)\) eventually, since
\(b(x)\asymp(\log x)^2\) while \(\log\beta(x)\asymp\log x\). Thus
\(e^{-b/2}\le\beta^{-c}\) eventually. Choose
\(c>\max(0,-q-z)\); the product rate \(-q-c\) is then less than \(z\),
so the strict-power definition gives the claimed rate. Taking \(z=-1\) and using the negative-rate consequence
above proves (4).
\end{proof}
